\section{Step 3 Form one scalar difference}\label{h263-sad-s03}
Use separate row-start pointers and a count of sixteen rows. At the start of a
row, copy those starts into moving sample cursors and initialize a sixteen-sample
column count. Load one unsigned byte from each cursor. Their subtraction lies
between minus 255 and plus 255, so its absolute value fits an ordinary general
register without a special minimum-integer case.

Accumulate the magnitude in a wide sum. The current input's first worked sample
is 56 and the reference sample is 200, giving magnitude 144. Loading 200 with a
signed byte load would instead produce minus 56, so the wrong result can arise
before the absolute-value operation. Interpret the samples correctly at the load.

\paragraph{Observe} Stop immediately after forming the first magnitude. It is
144. No input byte has been modified.
\Needspace{8\baselineskip}
\paragraph{Check} Which loaded values should be subtracted for this pair?
\begin{itemize}
  \item[A.] 56 and minus 56.
  \item[B.] Minus 200 and 56.
  \item[C.] 56 and 200.
\end{itemize}
